\begin{abstract}
Activation steering---adding a direction vector to a model's internal representations at inference time---is a simple and effective technique for controlling language model behavior. Nearly all prior work applies steering vectors to the residual stream, a $d_{\text{model}}$-dimensional space that aggregates information from all model components. We ask whether the MLP intermediate space, a $4\times$ wider post-activation representation where the model performs much of its computation, is a better or worse target for steering. We conduct a systematic comparison across three intervention points in GPT-2 Small: the residual stream (768-dim), the MLP intermediate space (3072-dim, post-GELU), and the MLP output (768-dim, after down-projection). We find that steering in the MLP intermediate space is possible but approximately $2.6\times$ less efficient than residual stream steering, as measured by sentiment probability shift per unit of KL divergence. Linear probes achieve lower accuracy in the MLP space (76\%) than the residual stream (86\%), contradicting the hypothesis that the wider space is more linearly separable. However, the MLP steering vector is remarkably sparse: just 2.8\% of neurons carry 50\% of the steering signal. This concentration makes MLP steering fragile---generations degrade to gibberish at lower intervention strengths. Our results suggest that the residual stream's advantage stems from its role as an integrative representation, while MLP neurons encode specific, less universal features that resist broad behavioral steering.
\end{abstract}
